\section*{Chapter \ref{ch:b1:competitive-industry} --- The Competitive Industry}
% Numbers recomputed in code/bachelor-1/12-competitive-industry/exercises/solutions.py (tested).

\begin{solution}{exo:b1:competitive-industry:sum}
$Q = 50 \times p/2 = 25p$.
\end{solution}

\begin{solution}{exo:b1:competitive-industry:entry}
$AC = 100/q + q$ is lowest at $q = 10$, where it is 20: the long-run price is
20. Demand is $1000 - 200 = 800$, so there are $800/10 = 80$ firms.
\end{solution}

\begin{solution}{exo:b1:competitive-industry:read}
United Kingdom: $181.9/35.8 = 5.08$. United States: $158.5/84.8 = 1.87$.
\end{solution}

\begin{solution}{exo:b1:competitive-industry:demand-rise}
In the short run, with the same firms, the price rises along the
short-run \omterm{def:b1:competitive-industry:industry-supply}{industry supply} and each firm makes a profit. Firms enter until
the price is back at the minimum of average cost; each produces as before,
so the number of firms rises by $20\%$.
\end{solution}

\begin{solution}{exo:b1:competitive-industry:increasing}
$p = 10 + (2000 - 50p)/100$ gives $1.5p = 30$: $p = 20$, $Q = 1000$, and a rent of
$\tfrac12 \times 1000 \times 10 = 5000$. With $Q = 2600 - 50p$: $1.5p = 36$, $p = 24$, $Q = 1400$ and the
rent is $\tfrac12 \times 1400 \times 14 = 9800$: the owners of the scarce input gain most of the
rise in demand.
\end{solution}

\begin{solution}{exo:b1:competitive-industry:quasi}
$4q = 24$ gives $q = 6$. \omterm{def:b1:competitive-industry:quasi-rent}{Quasi-rent} $24 \times 6 - 2 \times 36 = 72$; profit $72 - 50 = 22$.
\end{solution}

\begin{solution}{exo:b1:competitive-industry:licence}
$40000/0.05 = 800000$ dollars.
\end{solution}

\begin{solution}{exo:b1:competitive-industry:more-licences}
240 boats supply $Q = 60p$, so $60p = 3000 - 50p$ and $p = 27.27$; each boat lands
$6.82$ tonnes and earns $27.27 \times 6.82 - 50 - 2 \times 6.82^2 = 42.96$ with these rounded
values ($42.975$ unrounded). The licence is worth $42.975/0.04 = 1074.4$ thousand dalens,
$31.2\%$ less than 1,562.5: a $20\%$ increase in licences cuts their value by
almost a third.
\end{solution}

\begin{solution}{exo:b1:competitive-industry:growth}
United Kingdom: $5.08^{1/70} - 1 = 2.35\%$ a year. United States:
$1.87^{1/70} - 1 = 0.90\%$ a year.
\end{solution}

\begin{solution}{exo:b1:competitive-industry:subsidy}
A subsidy lowers each firm's average cost by $s$ at every output. In a
\omterm{def:b1:competitive-industry:cost-industry}{constant-cost industry} the long-run supply is flat at $\min AC$, which
falls by $s$: the price falls by $s$ and buyers gain all of it. If the
long-run supply rises because an input is scarce, the extra output raises
the price of that input; the long-run supply curve shifts down by $s$ but the
new price falls by less than $s$ (unless demand is perfectly elastic), and
the rise in the input's rent takes the rest.
\end{solution}

\begin{solution}{exo:b1:competitive-industry:expected}
With warning, the price of a medallion falls at the announcement to the
present value of ten years of rents: the owners at that date lose the
rest, and those who buy afterwards pay a price that already allows for the
end, so they lose nothing more. Without warning, the whole value is lost
on the day by whoever owns the medallions then.
\end{solution}

\begin{solution}{exo:b1:competitive-industry:refute}
If house prices rose as much where land is plentiful and building is easy
as where land is scarce, or if the share of land in house values did not
rise while prices did (so that construction costs drove the rise), the
statement would be refuted. Between places: prices in cities that can
expand freely against cities hemmed in by sea or rules. Over time: land
values from separate appraisals of land and buildings.
\end{solution}

\begin{solution}{pb:b1:competitive-industry:fishing}
\textbf{1.} $MC = 4q$ and $AC = 50/q + 2q$, lowest where $50/q^2 = 2$, at $q = 5$, where it is 20.

\textbf{2.} $q = p/4$ (above the price 20 in the long run).

\textbf{3.} 20 thousand dalens a tonne.

\textbf{4.} $3000 - 50 \times 20 = 2000$ tonnes.

\textbf{5.} $2000/5 = 400$ boats.

\textbf{6.} Still 20; the catch is 2,500 tonnes, from 500 boats.

\textbf{7.} $Q = 200 \times p/4 = 50p$.

\textbf{8.} $50p = 3000 - 50p$: $p = 30$, $Q = 1500$.

\textbf{9.} $q = 30/4 = 7.5$ tonnes; $AC = 50/7.5 + 15 = 21.67$.

\textbf{10.} $30 \times 7.5 - 50 - 2 \times 56.25 = 225 - 50 - 112.5 = 62.5$ thousand dalens a year.

\textbf{11.} $225 - 112.5 = 112.5$ thousand dalens.

\textbf{12.} Entry needs a licence, and the number of licences is fixed.

\textbf{13.} 62,500 dalens a year: the profit above the normal return already
counted in the fixed cost.

\textbf{14.} $62.5/0.04 = 1562.5$ thousand dalens: 1.5625 million.

\textbf{15.} $62.5/1562.5 = 4\%$: the normal return.

\textbf{16.} The fishers who received the licences when the cap was set: each
gained 1.5625 million dalens.

\textbf{17.} $62.5/0.025 = 2500$ thousand dalens.

\textbf{18.} Ten years of rents: $62.5 \times (1 - 1.04^{-10})/0.04 = 62.5 \times 8.1109 = 506.93$
thousand dalens.

\textbf{19.} $1562.5 - 506.93 = 1055.57$ thousand dalens.

\textbf{20.} The price falls to 20 and 400 boats fish.

\textbf{21.} From 1.5625 million to nothing: $100\%$.

\textbf{22.} Demand reaches zero at $p = 60$. With the cap,
$\tfrac12 \times 1500 \times (60 - 30) = 22500$; without it, $\tfrac12 \times 2000 \times 40 = 40000$ (thousands of
dalens a year).

\textbf{23.} Buyers gain 17,500 a year; licence owners lose $200 \times 62.5 = 12500$ a
year. The difference, 5,000, is what the cap cost: boats fishing above
their efficient size, $(21.67 - 20) \times 1500 = 2500$, and fish worth more than
its cost left uncaught, $\tfrac12 \times 500 \times 10 = 2500$.

\textbf{24.} Whoever holds the licences when the cap goes, including those
who bought them at 1.5625 million.

\textbf{25.} The capitalised value of one licence is 1.5625 million dalens;
it falls by $100\%$ when the cap goes.
\end{solution}
